\section{Related work}
\label{sec:related}

We organize the literature into four themes: extensions of the IDR framework for reproducibility analysis, Bayesian nonparametric copula models, Pitman--Yor mixture theory, and modern multiple-testing methodology relevant to the multi-replicate setting. Table~\ref{tab:related-work} summarizes the principal IDR-style competitors along the capability dimensions that matter for our method. Methods in the table should be read as direct baselines, methodological neighbors, or complementary approaches depending on their inferential target.

\begin{table}[H]
\centering
\caption{Comparison of IDR variants and related reproducibility methods along key capability dimensions. ``SC step-up'' refers to the Sun--Cai-type step-up rule applied to the local idr.}
\label{tab:related-work}
\small
\setlength{\tabcolsep}{3pt}
\begin{tabular}{lccccc}
\toprule
Method & Max $K$ & Copula family & Marginal model & FDR control & Theory \\
\midrule
Vanilla IDR \citep{LiBrownHuangBickel2011}        & 2          & Gaussian, fixed         & shared latent Gaussian & SC step-up & EM consistency \\
IDR2D \citep{Krismer2020IDR2D}                    & 2          & Gaussian, fixed         & shared latent Gaussian & SC step-up & --- \\
ChIP-R \citep{Newell2021ChIPR}                    & $\geq 3$   & none (rank product)     & rank-based             & p-value     & --- \\
eCV \citep{Iakovidis2024eCV}                      & $\geq 3$   & implicit (CV-based)     & shared                 & SC step-up & --- \\
nestedIDR \citep{Ranalli2026nestedIDR}            & $\geq 2$ per source & Gaussian, per-source & shared per source & SC step-up & --- \\
MaRR \citep{Philtron2018MaRR}                     & $\geq 2$   & none (rank-based)       & nonparametric          & frequentist & consistency \\
GMCM-AD \citep{KasaRajan2020GMCM}                 & $\geq 2$   & Gaussian mixture, fixed $K$ & shared latent Gaussian & --- & --- \\
NPMLE replicability \citep{Yi2024NPMLE}           & $\geq 2$   & none (density)          & nonparametric          & frequentist & consistency \\
BNP Archimedean copula \citep{PanNietoBarajasCraiu2024} & arbitrary & PY Archimedean, scalar param. & --- & --- & --- \\
\textbf{PY-IDR (this paper)}                      & arbitrary  & PY Gaussian + heavy-tail mixture & per-replicate Bernstein & SC plug-in step-up & contraction + mFDR \\
\bottomrule
\end{tabular}
\end{table}


\subsection{Extensions of IDR for reproducibility analysis}

The original IDR framework \citep{LiBrownHuangBickel2011} has been extended in several directions, both concurrently with and after the original publication. The eCV extension \citep{Iakovidis2024eCV} accommodates $K \geq 3$ replicates by replacing the bivariate Gaussian latent model with an aggregated coefficient-of-variation statistic, while retaining a single dependence-regime summary. The nestedIDR framework of \citet{Ranalli2026nestedIDR} introduces hierarchical replicability across multiple data sources via a nested mixture structure and is the most direct hierarchical IDR comparator in our planned study; implementation availability and command-line defaults are left as a coding-agent verification task. The maximum-rank-reproducibility method of \citet{Philtron2018MaRR} replaces the parametric Gaussian-copula mixture with a nonparametric ranking-based criterion, providing robustness to model misspecification at the cost of weaker per-feature posterior interpretability. IDR2D \citep{Krismer2020IDR2D} extends the framework to genomic interactions, and ChIP-R \citep{Newell2021ChIPR} extends to $K \geq 3$ ChIP-seq and ATAC-seq replicates via rank-product aggregation. \citet{KasaRajan2020GMCM} provide automatic-differentiation MLE for Gaussian mixture-copula models including the IDR special case. The robust NPMLE replicability method of \citet{Yi2024NPMLE} provides a frequentist alternative that estimates joint $p$-value densities without explicit copula dependence modeling. We benchmark against this collection of methods where software and input-format compatibility can be verified.

The replicability framework underlying these methods is reviewed by \citet{BogomolovHeller2023Review}, which together with \citet{BogomolovHeller2018} and \citet{Bogomolov2023PartialConjunction} provides modern theory for replicability across multiple studies. \citet{HellerYaacoby2014repfdr} introduce the closely related repfdr empirical-Bayes procedure for GWAS replicability, which informs our posterior decision rule.

\subsection{Bayesian nonparametric copula models}

Bayesian nonparametric copula modeling has a substantial recent literature. \citet{BurdaProkhorov2014} establish a copula-based factorization of Bayesian multivariate infinite-mixture models. \citet{MurrayDunsonCarinLucas2013} introduce Bayesian Gaussian copula factor models for mixed data, and \citet{SmithKhaled2012} treat Bayesian estimation of copula models with discrete margins. More recent contributions include the Dirichlet-process Student-$t$ copula mixture of \citet{LiuLi2025MixT} with explicit tail-dependence modeling, the conditional vine-copula formulation of \citet{DallaValleLeisenRossiniZhu2024}, the implicit copula variational inference framework of \citet{SmithLoaizaMayaNott2023}, the nonparametric copula models for multivariate, mixed, and missing data of \citet{FeldmanKowal2024NPCopula}, and the dependence-grouping mixture of \citet{ZhuangDiaoYi2022}.

The closest published methodological neighbor to our method is \citet{PanNietoBarajasCraiu2024}, which develops a Bayesian nonparametric mixture of Archimedean copulas with Pitman--Yor mixing. Our paper differs by (i) using Gaussian copulas with structured correlation matrices rather than Archimedean families with scalar dependence parameters; (ii) targeting the IDR reproducibility problem rather than general dependence modeling; and (iii) connecting posterior copula learning to local-idr ranking and Bayes-mFDR thresholding.

\subsection{Pitman--Yor mixture theory}

The Pitman--Yor process \citep{PitmanYor1997} extends the Dirichlet process to a two-parameter family with discount $\sigma \in [0,1)$ and concentration $\alpha > -\sigma$, with the Blackwell--MacQueen-type predictive characterization of \citet{Pitman1996BlackwellMacQueen}. Inference algorithms for PY mixtures are reviewed by \citet{IshwaranJames2001} and generalized by \citet{FavaroTeh2013} to normalized random measures with independent increments. Slice sampling for stick-breaking priors is developed by \citet{Walker2007SliceSampling} and \citet{KalliGriffinWalker2011}; the retrospective sampler of \citet{PapaspiliopoulosRoberts2008} provides a complementary computational strategy.

For posterior contraction rates, \citet{ShenTokdarGhosal2013} establish Hellinger contraction at the near-minimax rate $n^{-\beta/(2\beta+d)}(\log n)^t$ for DP location-scale mixtures of multivariate Gaussians on $\R^d$ under tail conditions. \citet{CanaleDeBlasi2017} generalize to broader covariance priors, and \citet{Scricciolo2014BAPY} provides adaptive rates in $L^p$-metrics specifically for Pitman--Yor and normalized inverse-Gaussian process kernel mixtures. Adaptive Bernstein-polynomial contraction rates for densities on $[0,1]$ are due to \citet{Rousseau2010BetaMix}, with the slower non-adaptive rate of \citet{Ghosal2001Bernstein} as the foundational result. Location-scale mixture contraction theory more generally is developed by \citet{KruijerRousseauVdV2010}.

Recent quantitative-mixing results for nonparametric mixture samplers are pioneered by \citet{Ascolani2024Mixing}, who give explicit polynomial mixing-time bounds for marginal Gibbs samplers in DPM models. Variational consistency results applicable to nonparametric models include \citet{AlquierRidgway2020}, \citet{YangPatiBhattacharya2020}, and \citet{ZhangGao2020}; the parametric Bernstein--von Mises result of \citet{WangBlei2019} applies most directly to truncated finite-dimensional models. Pitman--Yor applications in genomics include the hierarchical PY single-cell formulation of \citet{Bystrova2022HPY}.

\subsection{Modern multiple-testing methodology}

The Sun--Cai compound-decision framework \citep{SunCai2007,CaiSunWang2019} provides the theoretical underpinning for plug-in local-fdr procedures and is the basis of the IDR threshold rule. Recent developments include the optimal control-rate analysis of \citet{HellerRosset2021} and the empirical-Bayes spike-and-slab framework of \citet{CastilloRoquain2020}, which motivate our conditional rate-to-mFDR argument. Side-information-aware FDR procedures \citep{LeiFithian2018AdaPT,IgnatiadisHuber2021IHW} and $e$-value-based FDR control \citep{WangRamdas2022} are complementary rather than directly comparable, since they target different inferential objects.
